% !TEX program = xelatex

\documentclass[13pt]{extarticle}

% ---- Polices et encodage ----
\usepackage{fontspec}
\setmainfont{Verdana}
\setsansfont{Verdana}
\usepackage[french]{babel}

% ---- Mise en page ----
\usepackage[margin=2.5cm]{geometry}
\usepackage{setspace}
\onehalfspacing % interligne 1,5

% ---- Mathématiques ----
\usepackage{amsmath, amssymb}
\usepackage{mathpazo} % charge un package math compatible ; optionnel

% ---- Titres soulignés ----
\usepackage{titlesec}
\usepackage{ulem}

% Titres de parties (I, II, III) soulignés
% \uline placé dans le before-code (et non dans le format) pour éviter
% l'affichage en double du titre avec titlesec + ulem.
\titleformat{\section}
  {\normalfont\Large\bfseries}
  {\thesection}{1em}{\uline}
\titlespacing*{\section}{0pt}{1.5ex plus 1ex minus .2ex}{1.2ex plus .2ex}

% Titres de sous-parties soulignés
\titleformat{\subsection}
  {\normalfont\large\bfseries}
  {\thesubsection}{1em}{\uline}
\titlespacing*{\subsection}{0pt}{1.3ex plus 1ex minus .2ex}{1ex plus .2ex}

% ---- Commandes personnalisées pour souligner les types de blocs ----
\newcommand{\blocTitre}[1]{%
  \vspace{1em}%
  \noindent\underline{\textbf{#1 :}}\par%
  \vspace{0.3em}%
}

% ---- En-tête de page ----
\pagestyle{empty}

\begin{document}

% ===================== TITRE PRINCIPAL =====================
\begin{center}
  \Large\bfseries\underline{Séquence n\textsuperscript{o} 1 : Divisibilité et nombres premiers}
\end{center}

\vspace{1em}

% ===================== PARTIE I =====================
\section{Multiples et diviseurs}

\blocTitre{Définition}

Prenons $n$ et $k$ deux entiers. On dit que $n$ est un \textbf{multiple} de $k$ s'il existe un entier $y$ tel que :
\[
n = k \times y
\]

On dit alors aussi que $k$ est un \textbf{diviseur} de $n$, ou bien que $k$ \textbf{divise} $n$, ou bien encore que $n$ est \textbf{divisible} par $k$.

\blocTitre{Exemple}

\begin{itemize}
  \item 35 est un multiple de 5.
\end{itemize}

En effet, $35 = 5 \times 7$.

Il existe donc un entier tel que 5 multiplié par cet entier donne 35. Cet entier est 7.

Formulé autrement :
\begin{itemize}
  \item 5 est un diviseur de 35.
  \item 5 divise 35.
  \item 35 est divisible par 5.
\end{itemize}

% --- Page 2 ---

\blocTitre{Propriété}

Prenons $n$ et $k$ deux entiers.

$k$ divise $n$ si et seulement si $\dfrac{n}{k}$ est un entier.

\blocTitre{Exemples}

\textbf{16 est-il un multiple de 5 ?}

\[
\frac{16}{5} = 3{,}2
\]
et $3{,}2$ n'est pas un entier, donc 16 n'est pas un multiple de 5.

\textbf{119 est-il divisible par 7 ?}

\[
\frac{119}{7} = 17
\]
et 17 est un entier, donc 119 est divisible par 7.

% --- Page 3 ---

\blocTitre{Propriété : critères de divisibilité}

Prenons $n$ un entier.

\begin{enumerate}
  \item $n$ est divisible par \textbf{2} si et seulement si le chiffre des unités de $n$ est : 0, 2, 4, 6 ou 8.
  \item $n$ est divisible par \textbf{5} si et seulement si le chiffre des unités de $n$ est 0 ou 5.
  \item $n$ est divisible par \textbf{10} si et seulement si le chiffre des unités de $n$ est 0.
  \item $n$ est divisible par \textbf{3} si et seulement si la somme des chiffres permettant d'écrire $n$ est un multiple de 3.
\end{enumerate}

\blocTitre{Exemple}

235 :
\begin{itemize}
  \item 235 n'est pas divisible par 2 ; son chiffre des unités n'est pas 0, 2, 4, 6 ou 8.
  \item 235 est divisible par 5 ; son chiffre des unités est 5.
  \item 235 n'est pas divisible par 10 ; son chiffre des unités n'est pas zéro.
  \item 235 n'est pas divisible par 3 car $2 + 3 + 5 = 10$ et 10 n'est pas un multiple de 3.
\end{itemize}

% ===================== PARTIE II =====================
\section{Nombres premiers}

\blocTitre{Définition}

Un \textbf{nombre premier} est un entier différent de 1 et divisible seulement par lui-même et 1.

\blocTitre{Exemples}

2, 3, 5, 7, 11 sont des nombres premiers.

% --- Page 5 ---

\blocTitre{Théorème}

Tout nombre entier différent de 0 et 1 se décompose en produit de nombres premiers.

Cette décomposition est \textbf{unique}, à l'ordre des facteurs près.

\blocTitre{Exemple}

\begin{itemize}
  \item Décomposition de 66 :
\end{itemize}

\[
\begin{aligned}
66 &= 2 \times 33 \\
   &= 2 \times 3 \times 11
\end{aligned}
\]

On a décomposé 66 en produit de nombres premiers.

\begin{itemize}
  \item Décomposition de 83 :
\end{itemize}

\[
\begin{aligned}
83 &= 83
\end{aligned}
\]

83 est premier.

% --- Page 6 ---

\[
\begin{aligned}
\bullet\ 60 &= 2 \times 30 \\
            &= 2 \times 3 \times 10 \\
            &= 2 \times 3 \times 2 \times 5
\end{aligned}
\]

% ===================== PARTIE III =====================
\section{Fractions irréductibles}

\blocTitre{Définition}

Une fraction est \textbf{irréductible} si on ne peut pas la simplifier.

\blocTitre{Exemples}

\begin{itemize}
  \item $\dfrac{2}{3}$ est irréductible.
  \item Rendons $\dfrac{12}{18}$ irréductible.
\end{itemize}

\textbf{Méthodes :}

\medskip

\textbf{Méthode 1 :}
\[
\begin{aligned}
\frac{12}{18} &= \frac{2 \times 6}{2 \times 9} \\
              &= \frac{6}{9} \\
              &= \frac{2 \times 3}{3 \times 3} \\
              &= \frac{2}{3}
\end{aligned}
\]

\textbf{Méthode 2 :}
\[
\begin{aligned}
\frac{12}{18} &= \frac{6 \times 2}{6 \times 3} \\
              &= \frac{2}{3}
\end{aligned}
\]

\textbf{Méthode 3 :}
\[
\begin{aligned}
\frac{12}{18} &= \frac{2 \times 2 \times 3}{2 \times 3 \times 3} \\
              &= \frac{2}{3}
\end{aligned}
\]

\end{document}
